\chapter*{Afterword}
\addcontentsline{toc}{chapter}{Afterword}

We hope that this book has given you a solid foundation in Common Lisp\index{Common Lisp} and has inspired
you to explore this powerful language further. Common Lisp represents one of the oldest
continuously-used programming languages, yet it remains remarkably modern and relevant
for today's software development challenges.

\section*{The Power of Lisp}

The features that made Lisp powerful in the 1960s --- dynamic typing, automatic memory
management, interactive development, and code-as-data --- are now recognized as essential
features in modern programming languages. What sets Common Lisp apart is that it has
had these features from the beginning, refined over decades of practical use.

The ability to write code that writes code, through macros and metaprogramming, gives
Common Lisp a flexibility that few other languages can match. This power allows you to
create domain-specific languages embedded within CL, effectively raising the level of abstraction at which you program.

\section*{The Lisp Community}

One of the great strengths of Common Lisp is its community. Lisp programmers tend to be
thoughtful, experienced developers who value elegance and power in their tools. The community maintains a wealth of open-source libraries, documentation, and shared knowledge.

We encourage you to participate in the community:

\begin{itemize}
\item Ask questions on forums and mailing lists
\item Contribute to open-source Lisp projects
\item Write about your experiences with Lisp
\item Share your code and libraries
\item Help newcomers learn the language
\end{itemize}

\section*{Continuing Your Journey}

This book has covered the fundamentals, but there is much more to learn:

\begin{itemize}
\item Advanced macro techniques
\item The Common Lisp Object System\index{Common Lisp Object System} (CLOS\index{CLOS}) in depth
\item The condition system for sophisticated error handling
\item The reader macro system for extending CL syntax
\item Practical application development and deployment
\item Domain-specific languages built on CL
\item Performance optimization techniques
\end{itemize}

The books and resources listed in Appendix B will help you continue your learning. We
particularly recommend Paul Graham's ``On Lisp'' for advanced macro programming, and
Peter Norvig's ``Paradigms of Artificial Intelligence Programming'' for substantial examples
of CL in practice.

\section*{Common Lisp in Industry}

Common Lisp is used in production systems around the world, often in domains where
reliability, flexibility, and sophisticated problem-solving are critical:

\begin{itemize}
\item Knowledge-based engineering and CAD systems
\item Financial trading and risk analysis
\item Scheduling and optimization
\item Natural language processing
\item Bioinformatics
\item Web applications
\item Aerospace and defense
\end{itemize}

While Common Lisp may not have the market share of languages like Java or Python, it
has a dedicated user base and continues to be the language of choice for applications where
its unique strengths matter.

\section*{The Future of Common Lisp}

Common Lisp's strength lies in its stability. The ANSI standard ensures that code written
today will continue to work decades from now. At the same time, implementations continue
to improve, adding modern features like:

\begin{itemize}
\item Better integration with modern operating systems
\item Improved performance on current hardware
\item Support for contemporary protocols and formats
\item Enhanced development tools and IDEs
\item Growing libraries of open-source code
\end{itemize}

The emergence of the Quicklisp library manager has made it easier than ever to discover
and use Common Lisp libraries. The continued development of implementations like SBCL,
CCL, and Allegro CL\index{Allegro CL} ensures that Common Lisp remains a viable choice for serious software
development.

\section*{Final Thoughts}

Programming in Common Lisp is different from programming in most other languages. The
interactive, incremental development style can feel strange at first if you're used to the
edit-compile-link-run cycle of languages like C++ or Java. But once you become comfortable with the Lisp way of working, it's hard to go back.

The REPL (Read-Eval-Print Loop) isn't just a way to test small code snippets --- it's a
window into a living, breathing program that you're building and modifying in real time.
The ability to redefine functions, experiment with data structures, and test ideas immediately, all while your program continues to run, provides a development experience that is
difficult to match in other languages.

We encourage you to embrace the Lisp philosophy: start simple, build incrementally, and
let the language adapt to your problem rather than forcing your problem into the constraints
of the language.

\section*{Acknowledgments}

This book would not have been possible without the contributions of many people over the
years:

\begin{itemize}
\item The designers and implementers of Common Lisp, who created a language that has stood
the test of time
\item The Common Lisp community, who have shared their knowledge and code freely
\item Franz Inc., for their continued support of Common Lisp and their contributions to the
community
\item The many authors who have written books and tutorials about Lisp, making the language accessible to new generations of programmers
\end{itemize}

\section*{Closing}

Thank you for reading this book. We hope it has been helpful in your journey to learn
Common Lisp. Whether you're using CL for professional development, academic research,
or personal projects, we wish you success and enjoyment in your Lisp programming.

Remember: Lisp is not just a programming language --- it's a way of thinking about computation. Once you truly understand Lisp, you'll see programming differently, regardless of
what language you're using.

Happy Lisping!

\begin{flushright}
David J. Cooper, Jr.\\
February 14, 2011
\end{flushright}
